\documentclass{article}
\usepackage{amsmath,amssymb}
\usepackage{xurl}
\usepackage[hidelinks]{hyperref}
% Shared commands for notes in docs/paper-gaps/.

\DeclareUnicodeCharacter{2080}{\ifmmode{}_{0}\else\textsubscript{0}\fi}
\DeclareUnicodeCharacter{2081}{\ifmmode{}_{1}\else\textsubscript{1}\fi}
\DeclareUnicodeCharacter{2082}{\ifmmode{}_{2}\else\textsubscript{2}\fi}
\DeclareUnicodeCharacter{2099}{\ifmmode{}_{n}\else\textsubscript{n}\fi}

\newcommand{\Lean}{\textsc{Lean}}
\ExplSyntaxOn
\NewDocumentCommand{\leanid}{m}{
  \ifmmode
    \textsf{\detokenize{#1}}
  \else
    \group_begin:
    \urlstyle{sf}
    \tl_set:Nn \l_tmpa_tl {#1}
    \tl_replace_all:Nnn \l_tmpa_tl {\_} {_}
    \exp_args:NV \path \l_tmpa_tl
    \group_end:
  \fi
}
\ExplSyntaxOff
\providecommand{\tr}{\operatorname{tr}}
\newcommand{\tnleanIssueBase}{https://github.com/LionSR/TNLean/issues/}
\newcommand{\tnleanPrBase}{https://github.com/LionSR/TNLean/pull/}
\newcommand{\tnleanDocBase}{https://sirui-lu.com/TNLean/docs/}
\newcommand{\ghissue}[1]{\href{\tnleanIssueBase#1}{GitHub issue \##1}}
\newcommand{\ghpr}[1]{\href{\tnleanPrBase#1}{PR \##1}}
\newcommand{\doclink}[1]{\href{\tnleanDocBase#1.html}{\path{#1}}}

% Dirac notation and the identity operator, as in blueprint/src/macros/common.tex.
% \providecommand keeps note-local definitions authoritative.
\usepackage{dsfont}
\providecommand{\ket}[1]{|#1\rangle}
\providecommand{\bra}[1]{\langle #1|}
\providecommand{\braket}[2]{\langle #1 | #2 \rangle}
\providecommand{\ketbra}[2]{|#1\rangle\!\langle #2|}
\providecommand{\Id}{\mathds{1}}
\providecommand{\one}{\mathds{1}}
\providecommand{\C}{\mathbb{C}}
\providecommand{\MN}[1]{M_{#1}(\C)}
\providecommand{\MD}{M_D(\C)}

% External referencing for paper sources.  The build helper writes sanitized
% auxiliary files under build/paper-aux/ and excludes duplicated source labels.
\usepackage{xr}

% Import a generated paper auxiliary file with a caller-supplied label prefix.
% The candidate paths support builds from docs/paper-gaps/, the repository root,
% and build/paper-aux/ itself.
\newcommand{\importpaperaux}[2]{%
  \IfFileExists{../../build/paper-aux/#2.aux}{%
    \externaldocument[#1]{../../build/paper-aux/#2}%
  }{%
    \IfFileExists{build/paper-aux/#2.aux}{%
      \externaldocument[#1]{build/paper-aux/#2}%
    }{%
      \IfFileExists{#2.aux}{%
        \externaldocument[#1]{#2}%
      }{}%
    }%
  }%
}

% General external reference macros
\newcommand{\paperref}[2]{\ref{#1-#2}}
\newcommand{\papereqref}[2]{(\ref{#1-#2})}

% CPSV16 source (1606.00608 / MPDO-22-12-17-2.tex) external document and shortcuts
\importpaperaux{cpsv16-}{1606.00608/MPDO-22-12-17-2-tnlean}
\newcommand{\cpsvref}[1]{\paperref{cpsv16}{#1}}
\newcommand{\cpsveqref}[1]{\papereqref{cpsv16}{#1}}

% Verdict marker: \gapnote{<kind>}{<status>}. Kind is the policy
% classification (clarification, local-correction, scope-restriction,
% unfaithful, false-source, open-gap); status is open, wip (elimination
% underway), resolved, or historical. Machine-read by the site index;
% prints a verdict line.
\newcommand{\gapnote}[2]{\par\noindent\textbf{Verdict:} #1 (#2).\par}

\title{The unblocked parity reading of the Fundamental Theorem of matrix product unitaries}
\author{}
\date{2026-09-26}
\begin{document}
\maketitle
\gapnote{false-source}{open}

This verdict concerns the printed converse when its $U^{(N)}=V^{(N)}$
conclusion is read for the original, unblocked tensors. It does not classify
the corresponding blocked-family statement as false.

\section{The printed assertion}

Definition~SF of Cirac--P\'erez-Garc\'ia--Schuch--Verstraete defines the
standard form of the \emph{two-site block} $\mathcal U_2$ of a simple tensor
$\mathcal U$ \cite[lines~619--622]{Cirac2017MPU}. The ensuing diagram displays
$\mathcal U_2$ and $\mathcal V_2$. Nevertheless, the Fundamental Theorem of
MPU states that its local unitary relations, Equations~(SFuu) and~(SFvv),
are equivalent to $U^{(N)}=V^{(N)}$ for every $N$
\cite[lines~624--648]{Cirac2017MPU}. Read literally as an assertion about
the original, unblocked tensors, the converse from those local relations to
equality at odd lengths does not hold. The preceding diagram does, however,
support a coherent reading about the blocked tensors; this note does not
refute that blocked interpretation.

\section{A full-rank counterexample}

Take physical dimension $d=2$ and bond dimension $D=1$. With $i,j\in
    \{0,1\}$, define two one-site tensors by
\begin{align}
    \mathcal U^{ij}_{00} & =\delta_{ij},  &
    \mathcal V^{ij}_{00} & =-\delta_{ij}.
    \label{eq:mpu_standard_form_parity_letters}
\end{align}
Both have unitary periodic families. Their double-layer letters are the
one-by-one scalars $\delta_{ij}$: the two signs cancel under physical
adjunction. Thus the source's first simplicity contraction holds with
$a=b=(1)$, and its second, rank-one insertion identity is automatic in the
one-dimensional bond space. Both tensors are therefore simple. Their
normalized transfer matrices are the one-by-one identity; they have full bond
support and carry the canonical-form-II weight $\rho=(1)$.

We recall the source objects of Section~III of Ref.~\cite{Cirac2017MPU}, with
the leg orientation fixed in Ref.~\cite{gap:mpu_source_cut_orientation}. Write
$\mathcal U^i_{j,\alpha,\beta}$ for the entry of $\mathcal U^{ij}$ with left and
right auxiliary indices $\alpha,\beta$. The two \emph{source cuts} of
$\mathcal U$ are the flattenings
\begin{align}
    M_1(\mathcal U)_{(i,\beta),(\alpha,j)}
     & =\mathcal U^i_{j,\alpha,\beta}, &
    M_2(\mathcal U)_{(\alpha,i),(j,\beta)}
     & =\mathcal U^i_{j,\alpha,\beta}.
    \label{eq:mpu_standard_form_parity_cuts}
\end{align}
If $M_1$ has rank $k_1$ and $M_2$ has rank $k_2$, the source chooses compact
factorizations and right inverses
\begin{align}
    M_1           & =X_1Y_1,          &
    M_2           & =X_2Y_2,          &
    Y_1Z_1        & =I_{k_1},         &
    Y_2Z_2        & =I_{k_2},
    \label{eq:mpu_standard_form_parity_factors}
\end{align}
normalized by $X_1^\dagger(I_d\otimes\rho)X_1=I_{k_1}$ and $X_2^\dagger X_2=I_{k_2}$.
The gates are the contractions
\begin{align}
    u_{(l,r),(i_1,i_2)}
     & =\sum_\beta(Y_2)_{l,(i_1,\beta)}(Y_1)_{r,(\beta,i_2)}, &
    v_{(j_1,j_2),(r,l)}
     & =\sum_\alpha(X_1)_{(j_1,\alpha),r}(X_2)_{(\alpha,j_2),l}.
    \label{eq:mpu_standard_form_parity_gate_formulas}
\end{align}
Tildes denote the same objects for $\mathcal V$.

Neither tensor is zero, and both source cuts have rank two. In the natural physical coordinates,
\begin{align}
    M_1(\mathcal U) & =M_2(\mathcal U)=I_2,  &
    M_1(\mathcal V) & =M_2(\mathcal V)=-I_2.
\end{align}
Choose compact rank-two factorizations
\begin{align}
    X_1=X_2=Y_1=Y_2               & =I_2
                                  &        & \text{for }\mathcal U, \\
    \widetilde X_1=\widetilde X_2 & =-I_2, &
    \widetilde Y_1=\widetilde Y_2 & =I_2
                                  &        & \text{for }\mathcal V.
\end{align}
All singular values are one. The corresponding right inverses can be chosen as
$Z_1=Z_2=\widetilde Z_1=\widetilde Z_2=I_2$; the first weighted and second
ordinary $X$-normalizations both reduce to $I_2$ because $\rho=(1)$.
The source gates computed from these factors have, for both tensors, the entries
\begin{align}
    u_{(l,r),(p,q)} & =\delta_{lp}\delta_{rq}, &
    v_{(p,q),(r,l)} & =\delta_{pr}\delta_{ql}.
    \label{eq:mpu_standard_form_parity_gates}
\end{align}
Thus both gates are identity matrices in their displayed domain and codomain
orderings; in particular both are unitary. The two minus signs on the $X$
factors cancel in $v$. The printed local equations, in which $x,y,z$ are the
unitaries whose existence the Fundamental Theorem asserts, hold with
$x=y=-I_2$ and $z=1$: the product $x\otimes y$ acts trivially on
$u$, and the two minus signs on the open half-legs cancel in the shifted gate
$v$ of Equation~(SFvv). Equivalently, the two-site blocked tensors agree
letter by letter,
\begin{align}
    \mathcal V_2^{(i_1,i_2),(j_1,j_2)}
     & =(-\mathcal U^{i_1j_1})(-\mathcal U^{i_2j_2})
    =\mathcal U_2^{(i_1,i_2),(j_1,j_2)}.
\end{align}

On a periodic chain of length $N$, however,
\begin{align}
    V^{(N)} & =(-1)^N I_{2^N}, & U^{(N)} & =I_{2^N}.
    \label{eq:mpu_standard_form_parity_closed}
\end{align}
The families agree at every even length and disagree at every odd length;
for example, $V^{(3)}=-I_8\ne I_8=U^{(3)}$.

\section{Correct scope}

The local equations compare standard forms of $\mathcal U_2$ and
$\mathcal V_2$. Their converse can assert equality of the blocked periodic
families, equivalently $U^{(2N)}=V^{(2N)}$ for positive $N$, but cannot by
itself assert equality of the original families at every length. The forward
implication from equality of the original families to the local relations is
not affected. A converse about the original families needs an additional
condition detecting odd lengths; it is a corrected theorem, not the printed
source assertion. In particular, adding such a condition to a formal theorem
must not be presented as a proof of the unqualified source-labelled
equivalence. The virtual-unitary-gauge theorem obtained from equality at all
lengths $N>1$ concerns only the forward direction and remains valid.

The existing identities \leanid{MPOTensor.blockTensor_smul},
\leanid{MPOTensor.mpo_smul}, and \leanid{MPOTensor.mpo_idTensor} give a short
formal route to the two-site and periodic equalities in this example. The
formal parity witness records simplicity, equality of the two-site blocks,
and inequality at length three. The explicit factors and local $x,y,z$
relations above are the additional mathematical check of the printed local
antecedent; they are not asserted by block equality alone.

\bibliographystyle{alpha}
\bibliography{references}
\end{document}
